%%%%% KAPITEL: EINIGE HINWEISE %%%%%%%%%%%%%%%%%%%%%%%%%%%%%%%%%%%%%%%%%%%%%%%%%%%%%%
\chapter{Latex Tips und Tricks (To be deleted)}


Dies ist das zweite Kapitel.


\section{Fußnoten}
\label{sec:fussnoten}                 % Label für den Verweis auf diese Section


Erster Abschnitt.\footnote{Wir können eine Fußnote hinzufügen.} 



\section{Verweise auf Kapitel und Abschnitte}


Wir können einen Verweis auf Abschnitt \ref{sec:fussnoten} einfügen.



\section{Tabellen}


Wir können eine Tabelle einfügen, siehe Tabelle \ref{tab:beispieltab}.



\begin{table}[hbt]
\centering
\begin{tabular}{lcr}
\toprule
                 & \multicolumn{2}{c}{Verbundene Zellen} \\
\cmidrule{2-3}
   Erste Spalte  & Zweite Spalte  & Dritte Spalte \\
\midrule
   linksbündig   & zentriert      & rechtsbündig  \\
   ...           & ...            & ...           \\
\bottomrule
\end{tabular}
\caption{Beispiel-Tabelle}
\label{tab:beispieltab}
\end{table}



\section{Abbildungen}


Wir können auch Abbildungen einfügen, siehe Abbildung \ref{fig:beispielabb}.


\begin{figure}[hbt]
\begin{center}
\includegraphics[width=10cm]{images/hsba.png}   % für volle Breite verwendet man [width=\linewidth]
\caption{Beispiel-Abbildung} 
\label{fig:beispielabb}
\end{center}
\end{figure}



\section{Literaturverweise}


Wir können einen Artikel aus einem wissenschaftlichen Journal wie \textcite{Glove89} zitieren, ein Buch ist natürlich auch möglich (\cite{Goldb89}). Wir können auch Verweise auf Seitenzahlen einfügen: \textcite[][p. 78]{Goldb89}.
Außerdem können wir Sammelband-Beiträge wie \textcite{Steen01} und Konferenz-Beiträge wie \textcite{Merkl00} zitieren.
Es ist auch möglich, auf andere Quellen wie etwa Websites zu verweisen
(zum Beispiel \cite{HSBA18}).

Der Chicago-Stil für das Literaturverzeichnis ist in diesem Template bereits voreingestellt, ebenso wie der Auto-Jahr-Zitierstil.
Die Angaben zu den zitierten Quellen müssen in einer separaten Datei gespeichert sein (diese muss die Endung .bib haben).
Diese bib-Datei muss im Dokumentkopf angegeben werden (in diesem Template heißt sie Literatur.bib).
